\section{Uniform floors for commuting band products}

\begin{theorem}[Uniform floor for mean trees of commuting band products]
    \label{thm:replica_symmetric_product_mean_floor}
    \lean{TensorPower.ReplicaTransport.exists_meanTree_symBandMetric_prod_floor}
    \leanok
    \uses{lem:replica_symmetric_metric, def:mean_tree}
    Fix a finite tensor product $\mathcal V=\bigotimes_{v\in V}\C^{n_v}$,
    with $n_v>0$, and let $t>0$. There exists $C\ge0$ such that the
    following holds for every $K,k\ge0$, every label set $J$, and every
    family of partitions $\pi_{j,g}$ indexed by $j\in J$ and $0\le g<K$.
    Put $A_{j,g}=\widetilde A_k(\pi_{j,g})$, using the identity-extended
    band metric of Theorem~\ref{thm:replica_symmetric_band_floor}.
    Suppose that, for each $j$, the matrices $A_{j,g}$ commute at distinct
    positions $g$. Then every nonempty finite rooted weighted binary tree
    $\mathcal T$ with terminal labels in $J$ satisfies
    \begin{align}
        (k+2)^{-CK}I
        \le \mathcal T\bigl(A_{j,0}\cdots A_{j,K-1}:j\bigr).
    \end{align}
    The exponent $C$ is chosen before $K,k,J$, the partitions, the tree,
    and its interpolation parameters. The label set need not be finite.
    Empty site sets, $K=0$, $k=0$, repeated labels, and interpolation
    parameters at either endpoint are included; the empty product is $I$.
\end{theorem}

\begin{proof}
    \leanok
    \uses{thm:replica_symmetric_band_floor,
        thm:commuting_list_scalar_bound, thm:mean_tree_properties}
    Choose $C$ from Theorem~\ref{thm:replica_symmetric_band_floor} and
    set $b=(k+2)^{-C}>0$. The commuting-product bound gives
    $b^KI\le A_{j,0}\cdots A_{j,K-1}$ for each $j$.
    Since $b^K=(k+2)^{-CK}$, monotonicity and the constant-input identity
    for the mean tree give
    \begin{align}
        b^KI=\mathcal T(b^KI:j)
        \le\mathcal T\bigl(A_{j,0}\cdots A_{j,K-1}:j\bigr).
    \end{align}
\end{proof}

% Source: OpenAI2026AreaLaw, 07-comparators.tex, lines 603--621.
% These are scalar-floor consequences for the filtered norm.
\section{Filtered-norm bounds}

\begin{theorem}[Squared filtered norm under a scalar metric floor]
    \label{thm:scalar_floor_filtered_norm}
    \lean{Matrix.Transport.filteredNormSq_le_of_scalar_floor}
    \leanok
    Let $M$ be a positive definite complex matrix on a finite-dimensional
    space, which may have dimension zero. Suppose that $b>0$ and $bI\le M$.
    For any vector $\xi$, set $N^2=\|M^{-1/4}\xi\|^2$. Then
    \begin{align}
        N^2\le b^{-1/2}\|\xi\|^2.
    \end{align}
\end{theorem}

\begin{proof}
    \leanok
    \uses{thm:scalar_floor_nonpositive_power, thm:quadratic_form_order}
    Since $M$ is positive definite, spectral functional calculus gives
    $N^2=\operatorname{Re}(\xi^*M^{-1/2}\xi)$.
    Theorem~\ref{thm:scalar_floor_nonpositive_power} gives
    $M^{-1/2}\le b^{-1/2}I$. The quadratic-form order then gives the
    stated bound, using $\operatorname{Re}(\xi^*\xi)=\|\xi\|^2$.
\end{proof}

\begin{theorem}[Logarithmic filtered norm under a polynomial metric floor]
    \label{thm:polynomial_floor_filtered_log}
    \lean{Matrix.Transport.log_filteredNormSq_le_of_rpow_floor}
    \leanok
    Let $M$ be positive definite, let $C\ge0$ and $r\ge1$, and suppose
    that $r^{-C}I\le M$. For any vector $\xi$ with $\|\xi\|\le1$, set
    $N^2=\|M^{-1/4}\xi\|^2$. Then
    $\log N^2\le (C/2)\log r$.
    Here the real logarithm is assigned the value zero at zero.
\end{theorem}

\begin{proof}
    \leanok
    \uses{thm:scalar_floor_filtered_norm}
    If $\xi=0$, then $N=0$, and the right-hand side is nonnegative.
    Otherwise $N>0$, since $M^{-1/4}$ is invertible. The squared-norm
    bound gives $N^2\le(r^{-C})^{-1/2}\|\xi\|^2\le r^{C/2}$.
    Taking logarithms and using $\log(r^{C/2})=(C/2)\log r$ gives the result.
\end{proof}
